Personalizace znamená přizpůsobení obtížnosti a obsahu úloh výkonu studenta; AI sleduje, „jaký typ příkladů studentovi jde a jaký ne“, upravuje výběr/obtížnost a může identifikovat problematická slova a generovat cvičení \cite{kb_ai_101_tipu_pdf_04e4a115_page_8_1, kb_ai_101_tipu_pdf_04e4a115_page_7_1}.

Praktický postup pro pedagoga: \begin{enumerate}
\item Vyberte oblast/kompetenci pro personalizaci (např. algebra).
\item Nastavte parametry nástroje (rozsah obtížnosti, požadavek na postup) — nástroj vygeneruje příklady a úrovně \cite{kb_ai_101_tipu_pdf_04e4a115_page_8_1}.
\item Po odevzdání využijte přehled výkonu k cílenému opakování nebo krátkým intervenčním cvičením na detekované slabiny \cite{kb_ai_101_tipu_pdf_04e4a115_page_7_1}.
\end{enumerate}

Důsledky pro hodnocení: \begin{itemize}
\item Redesign úloh umožní varianty zadání a diferencované hodnoticí kritérium podle úrovní.
\item Požadování artefaktů procesu (postupy, skeny) pomáhá validovat autenticitu pokroku a slouží jako vstup pro další personalizaci \cite{kb_ai_101_tipu_pdf_04e4a115_page_8_1}.
\end{itemize}